\documentclass[aspectratio=169,dvipsnames,usenames]{beamer}
\usepackage{graphicx}
\usepackage{stmaryrd}
\usepackage[T1,T2A]{fontenc}
\usepackage[utf8]{inputenc}
\usepackage[english,russian]{babel}
\usepackage{amsmath}
\usepackage{amsfonts}
\usepackage{amssymb}
\usepackage{makeidx}
\usepackage{verbatim}
\usepackage{amsthm}
%\usepackage{bnf}
\usepackage{cmll}
\usepackage{tikz}
\usepackage{enumerate}
\usepackage{mathtext}
\usepackage{mathtools}
\usepackage{mathabx}
%\usepackage[left=2cm,right=2cm,top=2cm,bottom=2cm,bindingoffset=0cm]{geometry}
\usepackage{proof}
%\usepackage{paracol}
%\usepackage{enumitem}
\usepackage{xcolor}
\usepackage{colortbl}
%\usepackage{minted}
%\usepackage{hyperref}
\setbeamertemplate{navigation symbols}{}
\usetikzlibrary{graphs}
\usetikzlibrary{graphs.standard}
\usetikzlibrary{automata,positioning}
\usepackage{float}

\begin{document}

%\theoremstyle{dfn}
\newtheorem{dfn}{Определение}[section]
\newtheorem{nte}{Замечание}[section]

\newtheorem{axiom}{Аксиома}[section]
\newtheorem{thm}{Теорема}[section]
\newtheorem{lmm}[theorem]{Лемма}
\newtheorem{statement}{Утверждение}[section]
\newtheorem{oun_paragraph}{Пункт}[section]
\newtheorem{cons}{Следствие}[section]
\newtheorem*{exm}{Пример}

\newcommand{\comb}[1]{\operatorname{\mathcal{#1}}}
\newcommand{\func}[1]{\operatorname{#1}}
\newcommand{\reduction}[1]{{\color{OrangeRed}#1}}
\newcommand{\set}[1]{\left\{#1\right\}}

\def\from#1{\par \parbox{0.7\textwidth}{\par \hfill\raggedleft \it #1}} 

\begin{frame}{}
\begin{center}
{\LARGE Теория категорий\\Категорная логика\\Категорная абстрактная машина}
\end{center}
\end{frame}

\begin{frame}{Теория категорий}
\begin{dfn}[Категория $\mathcal{C}$]
\begin{itemize}
\item класс объектов $C$;
\item набор морфизмов $\text{Hom}_\mathcal{C}(A,B)$;
\item композиция: для $f: A \rightarrow B$ и $g: B \rightarrow C$ определена $g \circ f: A \rightarrow C$.
\item тождественный морфизм (единственный): $\text{id}_A: A \rightarrow A$
\end{itemize}

Причём:
\begin{itemize}
\item $h \circ (g \circ f) = (h \circ g) \circ f$;
\item $f \circ \text{id}_A = \text{id}_B \circ f = f$.
\end{itemize}
\end{dfn}

Формализовать теорию категорий можно по-разному.
\end{frame}

\begin{frame}{Примеры}
\begin{itemize}
\item Порядок.
\item Дискретная категория.
\item Set: морфизмы --- отображения
\item Top: морфизмы --- непрерывные отображения
\item Натуральные числа: единственный объект, морфизмы --- групповые операции.
\end{itemize}

Коммутативные диаграммы: прерывистая линия --- единственность (есть другие обозначения).
\end{frame}

\begin{frame}{Инициальные и терминальные объекты}
\end{frame}

\begin{frame}{Произведение и сумма объектов}
Сумма: пусть $Z, u: X\rightarrow Z, v: Y\rightarrow Z$. Тогда существует единственное $h: X+Y \rightarrow Z$,
что $h \circ i_X = u, h \circ i_Y = v$. Канонические вложения: $i_X$, $i_Y$.

\end{frame}

\begin{frame}{Произведение: канонические проекции.}

\end{frame}

\begin{frame}{Декартово-замкнутые категории}
Рассмотрим свойство: $(A^B)^C = A^{(B\cdot C)}$.

Его аналоги: карринг превращает $x: \langle A, B \rangle \rightarrow C$ в $x': A \rightarrow B \rightarrow C$.

Или теорема о дедукции: $A \with B \vdash C$ эквивалентно $A \vdash B \rightarrow C$.

Обобщение свойства: декартово-замкнутая категория.
\end{frame}

\begin{frame}{Экспоненциал}
\begin{dfn}
Экспоненциал $Z^Y$ --- объект и отображение оценки $\text{eval}: Z^Y \times Y \rightarrow Z$,
что для любого $X$ и $g : X \times Y \rightarrow Z$ найдётся $\lambda g$, что диаграмма коммутативна:
\end{dfn}

\includegraphics{ExponentialObject-01}
\end{frame}

\begin{frame}{Декартово-замкнутая категория}
\begin{dfn}Категория $\mathcal{C}$ --- декартово-замкнутая, если:
\begin{itemize}
\item в ней имеется терминальный объект;
\item в ней для любых $X$, $Y$ имеется $X \times Y$;
\item в ней для любых $Y$, $Z$ имеется $Z^Y$.
\end{itemize}
\end{dfn}

Пример: алгебра Гейтинга.
\end{frame}

\begin{frame}{Категорная логика}
Рассмотрим декартово-замкнутую категорию $\mathcal{C}$:

$\alpha \rightarrow \beta$ мы понимаем как $\alpha\vdash\beta$.

Заметим, что следующие правила/свойства мы можем доказать (здесь $\wedge$ --- произведение,
$\Leftarrow$ --- экспоненциал, $T$ --- терминальный объект):

\begin{tabular}{lll}
$\infer{A \rightarrow^{gf} C}{A \rightarrow^f B\quad B \rightarrow^g C}$\\
$A \rightarrow^{\circ_A} T$\\
$A \wedge B \rightarrow^{\pi_{A,B}} A$ & $A \wedge B \rightarrow^{\pi'_{A,B}} B$ & 
     $\infer{C\rightarrow^{\langle f,g\rangle}A \wedge B}{C \rightarrow^f A\quad C \rightarrow^g B}$\\

$(A \Leftarrow B) \wedge B \rightarrow^{\varepsilon_{A,B}} A$ & $\infer{C \rightarrow^{h^*} A\Leftarrow B}{C \wedge B \rightarrow^h A}$
\end{tabular}

\end{frame}

\begin{frame}{Декартово замкнутые категории, заданные уравнениями}
Давайте, вместо определений и картинок, зададим их системой уравнений. Для всех
$A,B,C$ и заданных на них стрелках $f,g,h$ выполнено следующее:
\vspace{0.5cm}

\begin{tabular}{ll}
$ f1_A = f,\quad 1_Bf = f,\quad (hg)f = h(gf) $ & (категории)\\
$ f = \circ_A$, если $f: A \rightarrow T$ & (терм. объект)\\\hline
$ \pi_{A,B}\langle f,g \rangle = f$, $\pi'_{A,B} \langle f,g \rangle = g$ & (произведение)\\
$ \langle \pi_{A,B} h, \pi'_{A,B} h \rangle = h$, если $f: C \rightarrow A$, $g: C \rightarrow B$, $h: C \rightarrow A \wedge B$\\\hline
$ \varepsilon_{A,B} \langle h^* \pi_{C,B}, \pi'_{C,B} \rangle = h$, если $h: C \wedge B \rightarrow A$, $k: C\rightarrow A \Leftarrow B$ & (дек.-замкн.)\\
$ (\varepsilon_{A,B} \langle k \pi_{C,B}, \pi'_{C,B} \rangle)^* = k$
\end{tabular}
\end{frame}

\begin{frame}{Комбинаторные уравнения}
Пусть $M$, $N$ --- какие-то лямбда-выражения.

Рассмотрим комбинаторы $S$, $\Lambda$, $\text{'}$:
$\llbracket M N \rrbracket = S(\llbracket M \rrbracket, \llbracket N \rrbracket)$, $\llbracket \lambda . M\rrbracket = \Lambda (\llbracket M \rrbracket)$,
$\llbracket c \rrbracket = \text{'}c$\vspace{0.5cm}

Добавим комбинатор пары: $\llbracket (M, N) \rrbracket = \langle \llbracket M \rrbracket, \llbracket N \rrbracket \rangle$,
и $Fst$, $Snd$, с очевидным смыслом.


\begin{tabular}{ll}
$(ass)$ & $(x \cdot y) z = x(yz)$\\
$(fst)$ & $Fst(x,y) = x$\\
$(snd)$ & $Snd(x,y) = y$\\
$(dpair)$ & $\langle x, y \rangle z = (xz, yz)$\\
$(ac)$ &  $App(\Lambda(x)y,z) = x (y,z)$\\
$(quote)$ & $('x)y = x$
\end{tabular}
\end{frame}

\begin{frame}{Категорная абстрактная машина}
Память: (терм) (код) (стэк)

Операции:


\begin{tabular}{llll}
push & < & положить терм на стек \\
swap & , & обменять терм с вершиной стека\\
cons & > & построить пару из терма и вершины стека\\
fst & $Fst$ & терм: $(x,y) \rightarrow x$\\
snd & $Snd$ & терм: $(x,y) \rightarrow y$\\
cur & Л & код: $(\textbf{cur } C);CC \rightarrow CC$, терм: $s \rightarrow C:s$\\
app & App & код: $\textbf{app};CC \rightarrow C;CC$, терм: $(CC:s,t) \rightarrow (s,t)$\\
quote & ' & код: $(\textbf{quote }c);CC \rightarrow c;CC$, терм: $s \rightarrow c$ 
\end{tabular}
\end{frame}

\end{document}
